\chapter{Hybrid compressed and addressable memory}

A persistent agent needs more than a larger context window. It needs explicit answers to four questions: what state is transient, what information becomes durable, what may be compressed, and what must remain addressable with provenance.

\section{Four memory classes}
We distinguish (1) engine-native transient state, (2) immutable durable events, (3) compressed summaries, and (4) addressable indexed records. These classes can refer to one another but are not interchangeable.
\Plate{01-memory-classes}{Evolutor memory sits above native DOGMA/Hermon execution state.}{evod40:classes}

DOGMA state may compact history into a recurrent/native representation. Hermon KV preserves token-addressable attention state. Neither is automatically an auditable long-term memory record.

\section{Durable record}
A durable record is
\[
R=(id,tenant,kind,payload,sources,t_{create},t_{expire},artifact).
\]
Identity and provenance survive retrieval and compression. Tenant and policy scope are part of semantics, not metadata decoration.

\section{Compression}
A compression operator
\[
C:\{R_1,\ldots,R_n\}\rightarrow S
\]
is generally non-injective. Multiple source sets can yield the same summary; a summary therefore cannot be treated as equivalent to its source events.
\Plate{02-compression}{Compression is lossy but provenance-preserving.}{evod40:compression}

The reference stores source IDs explicitly. A production compressor should also record model/version, prompt/policy, timestamp and quality checks.

\section{Addressability}
An addressable memory supports a query
\[
Q(q,\mathcal M,k)\rightarrow[(id_i,score_i)]_{i=1}^k.
\]
The index may be lexical, vector, graph, temporal or hybrid. The semantic requirement is that retrieved payloads can be traced to stable records.

\Plate{03-retrieval}{Retrieval returns ranked records, not anonymous context text.}{evod40:retrieval}

\section{Conditional retrieval}
Retrieval has cost and can introduce irrelevant or adversarial context. It should be invoked when a plan requires external/durable evidence rather than mechanically before every model step. Adaptive retrieval work such as Self-RAG illustrates this broader principle.

\section{Working set}
The active working set is a bounded selection of durable records promoted into the current execution context. Promotion does not transfer authority: a retrieved tool instruction remains untrusted data unless policy independently authorizes it.

\section{Rehydration}
Compressed memory can be expanded by following provenance links:
\[
H(S,B)\rightarrow \{R_{i_1},\ldots,R_{i_m}\},\quad m\le B.
\]
\Plate{04-rehydrate}{Rehydration is a bounded provenance traversal, not inversion of a lossy summary.}{evod40:rehydrate}

\section{Retention and forgetting}
Retention is a policy over time, relevance, legal requirements, user intent, storage budget and evidence value. Expiry makes a record unavailable after a declared time. Explicit forgetting creates a tombstone so stale indexes cannot silently resurrect it.
\Plate{05-lifecycle}{Deletion semantics must dominate indexes and caches.}{evod40:lifecycle}

\section{Provenance}
Every derived record carries source edges. This creates a memory DAG. Provenance supports auditing and targeted invalidation: if a source is deleted or corrected, dependent summaries can be marked stale and recomputed.

\section{Consistency}
The teaching oracle is single-process and strongly consistent. Distributed implementations must define whether index updates, object storage and tombstones are atomic, eventually consistent or transactionally coordinated. ``Memory'' without a consistency contract is underspecified.

\section{Memory safety firewall}
\Plate{06-firewall}{Memory content is evidence/context, never execution authority by itself.}{evod40:firewall}
Retrieved text can contain stale facts, prompt injection or malicious tool instructions. The planner may use it as data; the policy layer independently decides whether an action is permitted.

\section{Information accounting}
Let source information be $X$ and compressed representation $Z=C(X)$. Unless $C$ is injective, there is no general inverse $C^{-1}$. A useful system therefore measures compression ratio separately from task retention:
\[
\rho=\frac{|Z|}{|X|},\qquad
U=\operatorname{quality}(task\mid Z).
\]
Low $\rho$ is not success if $U$ collapses.

\section{Native state versus durable memory}
A native engine checkpoint can be useful for exact continuation but poor for semantic retrieval. A text/event record can be useful for retrieval but insufficient to restore numerical engine state. Evolutor may retain both, linked by provenance.

\section{MemGPT comparison}
MemGPT frames extended context as movement through a memory hierarchy. Evolutor adopts the explicit-tier lesson but adds a stronger separation between native engine state, durable event evidence, lossy summaries and addressable records because it must orchestrate two engine families.

\section{Executable reference}
\Listing{reference}
The reference implements immutable event records, lossy summaries with source IDs, deterministic lexical retrieval, TTL expiry, tombstone forgetting and bounded source rehydration.

The lexical ranker is intentionally simple. The chapter tests memory semantics, not retrieval SOTA.

\section{Failure modes}
Failures include cross-tenant retrieval, summary without provenance, deleted records surviving an index, unbounded rehydration, treating a summary as exact history, stale derived records, retrieval-as-authority and accidental serialization of secrets.

\section{Twelve worked exercises}
\begin{enumerate}
\item Is a KV cache durable memory? \textbf{Not by default; it is native transient execution state.}
\item Is compression invertible? \textbf{Only if the chosen transform is injective/lossless.}
\item Why source IDs? \textbf{Audit, rehydration and invalidation.}
\item Why tombstones? \textbf{To prevent stale indexes/caches from resurrecting deleted records.}
\item Retrieval result trusted? \textbf{No.}
\item What is promotion? \textbf{Selecting durable records into bounded working context.}
\item What does TTL express? \textbf{A declared availability/retention horizon.}
\item Why separate compression ratio and utility? \textbf{Smaller memory can destroy task-relevant information.}
\item Can an event restore DOGMA state? \textbf{Not unless a declared decoder/checkpoint relation exists.}
\item Can a DOGMA checkpoint answer semantic search directly? \textbf{Not necessarily.}
\item Distributed memory needs what extra contract? \textbf{Consistency, replication, deletion and conflict semantics.}
\item Release rubric? \textbf{Class separation; provenance; isolation; retention; tombstones; bounded retrieval/rehydration; injection firewall; utility tests.}
\end{enumerate}

\section{Bridge}
Chapter 41 turns the Chapter 39 runtime and Chapter 40 memory operators into a typed intermediate representation. Planning becomes compilation: validate effects and capabilities, optimize only semantics-preserving regions, allocate resources, then lower each operation to its legal backend.
